\cotefile{157-5}
\scanpage{149}
\mspage{65}
% archive sheet number: 148

\unsure{dont|donc} il devient évident à présent qu'il faut le faire de toutes façons. Je vais quand même \unsure{énoncer|résumer} le positif de mes perplexités :

\uline{Proposition-per\deleted{\unsure{p}}plexité} 4.3.1. Soient $M,N$ deux catégories pseudo-topos, d'où $M\boxtimes N$ comme dans 4.1. Soient $S\subset M$, $T\subset N$ deux petites sous-catégories strictement génératrices, et considérons l'inclusion pleinement fidèle correspondante

\diagram{\[
\begin{tikzcd}[ampersand replacement=\&,column sep=large]
\underset{\substack{\shortparallel\,\mathrm{def}\\\underline{\unsure{\mathrm{Hom}|\mathrm{Hum}}}^{!!}(M^{\unsure{o|0}}\times\unsure{N|N^o},\mathrm{Ens})}}{M\boxtimes N} \arrow[r,hook,"\alpha"'] \& S^\wedge\boxtimes T^\wedge \approxeq (S\times T)^\wedge \arrow[l,dashed,bend right=15,"{\varphi?}"']
\end{tikzcd}.
\tag*{(4.3.1.1)}
\]}

\inserted{Considérons} les conditions \unsure{suivantes|suiv.} \deleted{sont équivalentes}

a) $M\boxtimes N$ est un pseudo-topos

b) Pour $x\in\mathrm{Ob}\,M$, $y\in\mathrm{Ob}\,N$, $x\boxtimes y\in\mathrm{Ob}\,M\boxtimes N$ \unsure{existe}, i.e. le foncteur (commutant aux petites $\underleftarrow{\lim}$)

\[
F\longrightarrow F(x,y),\qquad M\deleted{\unsure{\times}}\boxtimes N\longrightarrow\mathrm{Ens}
\tag*{(4.3.1.2)}
\]

est représentable.

b$'$) Comme b), avec $x\in\mathrm{Ob}\,S$, $y\in\mathrm{Ob}\,T$.

c) Le foncteur $\alpha$ \deleted{\unsure{admet un adj.}} \inserted{(lequel commute aux petites $\underleftarrow{\lim}$)} admet un adjoint à gauche $\varphi:(S\times T)^\wedge\longrightarrow M\boxtimes N$.

On a alors les implications

\deleted{Démonstration}

\[
b\Longrightarrow b'\Longleftrightarrow c\mathrel{\deleted{\unsure{<}}}\Longrightarrow a.
\tag*{(4.3.1.3)}
\]

De plus, si $c$ est satisfait, on a \deleted{\unsure{précis|préci}} : pour $s\in S^\wedge$, $t\in T^\wedge$, \deleted{$\unsure{\rho}(s,t)$ par sa \unsure{définition}} considérons $(s,t)$ \deleted{$\unsure{\in}\, S^\wedge\boxtimes T^\wedge\simeq(S\times T)^\wedge$}\inserted{$\in(S\times T)\subset(S\times T)^\wedge\approxeq S^\wedge\boxtimes T^\wedge$}, \deleted{qui} s'explicite \deleted{\unsure{au}} comme

\margin{\diagram{\begin{tikzcd}[ampersand replacement=\&,column sep=small,row sep=small]
S\times T\arrow[r,"\mathrm{pr}_2"]\arrow[d,"\mathrm{pr}_1"] \& T\\
S \&
\end{tikzcd}}}

\[
\left\{\begin{aligned}
(s\deleted{\unsure{\times}},t)&=\mathrm{pr}_1^{\unsure{*|\alpha}}(s)\times\mathrm{pr}_2^{*}(t)
\deleted{\unsure{=(s,t)}}\inserted{\simeq s\overset{S^\wedge,T^\wedge}{\boxtimes}t}\in\mathrm{Ob}\,(S\times T)^\wedge\\
\text{i.e.}\qquad(s\deleted{\unsure{\times}},t)(s',t')&=s(s')\times t(t')=\mathrm{Hom}_{S}(s',s)\times\mathrm{Hom}_{T}(t',\unsure{t})\\
&=\mathrm{Hom}_{S\unsure{\times}T}((s',t'),(s,t)).
\end{aligned}\right.
\tag*{(4.3.1.4)}
\]

